% Standalone section. Requires amsmath, amssymb, amsthm and standard theorem environments.
% Source reviewed: ProveIt commit a0a90ef31877f98be437191c48b46f02d5456867.
\section{Cyclic prime symmetries and integral torsion}
\label{cpc:sec:main}

The incoming integral-distribution reports \cite{integraljets} determine reflection torsion
and explicitly propose a subgroup of prime order as a next target.
We solve that problem for an action which is trivial on one primary
factor and has no nonzero fixed point on the complementary factor,
assuming the fixed factor has semisimple unit-group action modulo the
symmetry prime. The result includes arbitrary integral weights and all
one-variable integral jets. Its obstruction is a finite multiplicative
compatibility condition, rather than a condition on each weight in
isolation.

The finite resolution and double-complex method have classical
antecedents in Kubert, Anderson, and Ouyang
\cite{cpc:Kubert1979,cpc:Ouyang2001,cpc:Ouyang2002}.
In particular, Ouyang's universal norm distributions already provide a
much broader setting for weighted distribution relations and their
cohomology. The calculation below is an explicit continuation for the
specified finite grid and a diagonal cyclic subgroup. We claim no global
priority for its ordinary-weight specialization or for the resolution
method.

\subsection{The formal module and the hypotheses}

For a commutative ring $B$, let $X_q=(1/q)\mathbb Z/\mathbb Z$ and
choose weights $a_p\in B$, one for each prime $p\mid q$. Define
\begin{equation}\label{cpc:eq:distribution}
 D_q(B,\mathbf a)=
 \frac{B[X_q]}{
 \left\langle\displaystyle\sum_{py=x}e_y-a_pe_x:
                    p\mid q,\ x\in X_{q/p}\right\rangle}.
\end{equation}
Every root sum is taken in $X_q$. The symbol $e_0$ is retained.
The integral normal-form theorem from the incoming reports, recalled
below with its proof mechanism, gives a $B$-free module of rank
$\varphi(q)$ for every specialization.

Let $\ell$ be a prime and let $u\in(\mathbb Z/q\mathbb Z)^\times$
have order $\ell$. Multiplication by $u$ induces
$T(e_x)=e_{ux}$. We impose the following explicit hypotheses:
\begin{equation}\label{cpc:eq:hypotheses}
 q=md,\qquad (m,d)=1,\qquad d>1,\qquad
 u\equiv1\pmod m,\qquad (u-1,d)=1,\qquad
 \ell\nmid\varphi(m).
\end{equation}
Here $\varphi(1)=1$. Necessarily $m=(u-1,q)$, so the factorization,
when it exists, is determined by $u$. The $T$-fixed grid is precisely
$X_m$. Write
\begin{equation}\label{cpc:eq:parameters}
 r=\omega(d),\qquad
 H=\langle p\bmod m:p\mid d\rangle
       \ \le (\mathbb Z/m\mathbb Z)^\times.
\end{equation}
The group modulo one is the trivial group.

For integer weights, call the constants \emph{compatible} when
\begin{equation}\label{cpc:eq:compatibility}
 p\bmod m\longmapsto \overline a_p\in\mathbb F_\ell^\times
 \quad(p\mid d)
\end{equation}
extends to a homomorphism $H\to\mathbb F_\ell^\times$.
In particular, all these residues must be nonzero. Equivalently,
every relation $\prod_{p\mid d}p^{n_p}\equiv1\pmod m$, with
$n_p\in\mathbb Z$, must satisfy
$\prod_{p\mid d}\overline a_p^{\,n_p}=1$.
Set
\begin{equation}\label{cpc:eq:multiplicity}
 c(\mathbf a)=
 \begin{cases}
 \varphi(m)/|H|,&\text{if \eqref{cpc:eq:compatibility} holds},\\
 0,&\text{otherwise}.
 \end{cases}
\end{equation}

\begin{theorem}[Integral cyclic-prime coinvariants]
\label{cpc:thm:scalar}
Assume \eqref{cpc:eq:hypotheses} and let all weights be integers.
Then
\begin{equation}\label{cpc:eq:scalar}
 \boxed{\displaystyle
 D_q(\mathbb Z,\mathbf a)/(T-1)D_q(\mathbb Z,\mathbf a)
 \simeq \mathbb Z^{\varphi(q)/\ell}
       \oplus (\mathbb Z/\ell\mathbb Z)^{\,c(\mathbf a)2^{r-1}}.}
\end{equation}
Consequently every nonunit nonzero Smith invariant of the complete
distribution-plus-symmetry presentation equals $\ell$, with the
multiplicity in \eqref{cpc:eq:scalar}. The answer is independent of
all weights $a_p$ with $p\mid m$.
\end{theorem}

\begin{corollary}[No nonzero fixed grid point]
\label{cpc:cor:fixedfree}
If $u$ has prime order $\ell$ and $(u-1,q)=1$, then
\[
 D_q/(T-1)D_q\simeq
 \mathbb Z^{\varphi(q)/\ell}\oplus
 \begin{cases}
 (\mathbb Z/\ell\mathbb Z)^{2^{\omega(q)-1}},
       &a_p\equiv1\pmod\ell\ \text{for every }p\mid q,\\
 0,&\text{otherwise}.
 \end{cases}
\]
For $\ell=2$ and odd $q$, this recovers the unanchored reflection
coinvariant formula. For $\ell>2$ it gives odd torsion.
\end{corollary}

The theorem describes a formal quotient. The equation $T=1$ is an
additional symmetry relation. Ordinary Hurwitz-zeta or polylogarithm
evaluation is not asserted to satisfy it without an appropriate
invariant projection. In particular, the finite torsion here is not
torsion of a complex-valued period and proves no numerical independence
statement.

\subsection{The finite resolution and its fixed part}

We record the algebra needed for a proof, including the role of the
prime-order assumption. Order the primes dividing $q$, and for a
subset $S$ put $p_S=\prod_{p\in S}p$. The weighted distribution
resolution has terms
\begin{equation}\label{cpc:eq:resolution}
 C_j=\bigoplus_{|S|=j}B[X_{q/p_S}],
 \qquad
 \partial[S,x]=\sum_{p\in S}(-1)^{\#\{v\in S:v<p\}}
 \left(\sum_{py=x}[S\setminus\{p\},y]
            -a_p[S\setminus\{p\},x]\right).
\end{equation}
Its augmentation is the quotient in \eqref{cpc:eq:distribution}.
The transfers and inclusions in different prime directions commute,
which gives $\partial^2=0$.

\begin{lemma}[Exactness over every coefficient ring]
\label{cpc:lem:resolution}
The augmentation of \eqref{cpc:eq:resolution} is exact. Its degree-zero
homology is $B$-free of rank $\varphi(q)$, and the augmented complex
is split exact as a complex of $B$-modules.
\end{lemma}
\begin{proof}
Here is the finite filtration argument, to make clear that no division
by $\ell$ or a prime weight is used. Filter $[S,x]$ by the positive
integer $p_S\operatorname{den}(x)$, which divides $q$. The leading
part of a prime differential retains only the preimages whose
denominator is multiplied by that prime; its weight term lies lower
in the filtration. At a fixed augmented denominator, additive primary
coordinates identify the leading complex with a tensor product of
maps
\[
 B[U_{p^{a-1}}]\longrightarrow B[U_{p^a}],
 \qquad e_b\longmapsto\sum_{v\mapsto b}e_v.
\]
At $a=1$, $U_1$ is a singleton and the sum has $p-1$ terms.
At $a>1$, each fiber has $p$ terms. In making this identification,
replace the $v$-primary coordinate in the $S$-summand by
$(\prod_{p\in S,\,p\ne v}p)^{-1}$ times that coordinate; this
removes the unit twists in the other primary components of a root.

Each displayed map is a split injection: its fibers have disjoint
support, and one coefficient-one pivot can be selected in each fiber.
Its cokernel is free of rank
$\varphi(p^a)-\varphi(p^{a-1})$. The tensor product is therefore
acyclic in positive degrees and has free degree-zero homology.
Lifting a leading boundary and subtracting it lowers the finite
filtration. This proves positive-degree exactness. The same lifting
in degree one proves that the image into degree zero is strict for
the filtration, so the chosen free complements lift to a basis of
the quotient. Their total number is
\[
 \prod_{p^e\parallel q}
 \left(1+\sum_{a=1}^{e}
      (\varphi(p^a)-\varphi(p^{a-1}))\right)=\varphi(q).
\]
Finally a surjection onto the free quotient splits; its kernel is
projective. Inducting upward splits the entire augmented complex.
\end{proof}

Let $C_j^{\rm mov}$ be generated by points which are not fixed by
$T$. It is a subcomplex: if $y$ is fixed, then $py$ is fixed, so a
root of a nonfixed point cannot be fixed. Since $T$ has prime order,
every nonfixed orbit has exactly $\ell$ elements, and each term of
$C_\bullet^{\rm mov}$ is a sum of regular cyclic permutation modules.
Put $K_\bullet=C_\bullet/C_\bullet^{\rm mov}$.

The fixed symbols in the summand indexed by $S$ form the grid
\[
 X_{m/p_{S\cap P(m)}}.
\]
For $p\mid m$, its fixed differential is the usual weighted
distribution differential. For $p\mid d$, exactly one root remains
fixed, namely $p^{-1}x$ in that grid. Thus this differential is
\begin{equation}\label{cpc:eq:fixedoperator}
 V_p-a_p,\qquad V_p(e_x)=e_{p^{-1}x}.
\end{equation}
All these operators commute. Applying Lemma~\ref{cpc:lem:resolution}
in the prime directions dividing $m$ shows that $K_\bullet$ is
quasi-isomorphic to the operator Koszul complex
\begin{equation}\label{cpc:eq:operatorKoszul}
 K\bigl(V_p-a_p:p\mid d;,D_m(B,\mathbf a|_m)\bigr).
\end{equation}
The notation means the tensor construction with one exterior
generator for each $p\mid d$, differential given by the indicated
commuting endomorphisms. The comparison follows by taking homology
first in the bounded $m$-directions; only degree zero survives.

\subsection{A cohomology formula which controls all torsion}

For a free abelian group $M$ with $T^\ell=1$, write
$\mathcal N=1+T+\cdots+T^{\ell-1}$ and define
\[
 \widehat H^0(M)=\ker(T-1)/\mathcal N M,
 \qquad
 \widehat H^1(M)=\ker\mathcal N/(T-1)M,
\]
extended periodically with period two. These are the cohomology
groups of the complex with differential $T-1$ leaving even degree
and $\mathcal N$ leaving odd degree.

\begin{lemma}[Torsion and the fixed complex]
\label{cpc:lem:tate}
For $B$ finite free over $\mathbb Z$ and $D=D_q(B,\mathbf a)$,
\begin{equation}\label{cpc:eq:torsiontate}
 \operatorname{Tor}_{\mathbb Z}(D/(T-1)D)
       =\widehat H^1(D).
\end{equation}
This group is killed by $\ell$. Under
\eqref{cpc:eq:hypotheses}, for $\epsilon\in\{0,1\}$ one has
\begin{equation}\label{cpc:eq:tateKoszul}
 \widehat H^\epsilon(D)\simeq
 \bigoplus_{j\equiv\epsilon\ ({\rm mod}\ 2)}
 H_j\!\left(
 K(V_p-\overline a_p:p\mid d;,
        D_m(B/\ell B,\overline{\mathbf a}|_m))\right).
\end{equation}
The isomorphism is one of $B$-modules; the right side is annihilated
by $\ell$.
\end{lemma}
\begin{proof}
If $kv=(T-1)w$ for a nonzero integer $k$, apply $\mathcal N$.
Freeness over $\mathbb Z$ gives $\mathcal N v=0$. Conversely,
if $\mathcal N v=0$, then $\ell v$ belongs to $(T-1)M$ because
$\ell-\mathcal N$ is divisible by $T-1$ in
$\mathbb Z[T]/(T^\ell-1)$. This proves both
\eqref{cpc:eq:torsiontate} and the annihilation assertion.

Apply the periodic cohomology complex to the finite resolution
\eqref{cpc:eq:resolution}. Taking horizontal homology replaces it
by $D$ in degree zero. A regular orbit has an exact periodic complex:
the image of $T-1$ is the coefficient-sum-zero subgroup, and the image
of $\mathcal N$ is the constant-vector subgroup. Hence removing
$C_\bullet^{\rm mov}$ does not change total cohomology. Both
comparisons are justified by finite horizontal width; each total
degree contains only finitely many summands.

On $K_\bullet$ the action is trivial, so the vertical complex
alternates between zero and multiplication by $\ell$. Since each
term of $K_\bullet$ is $\mathbb Z$-free, each adjacent
multiplication-by-$\ell$ block can be replaced by reduction modulo
$\ell$. The horizontal degree $j$ shifts total degree by $-j$,
which gives the parity in \eqref{cpc:eq:tateKoszul}. Finally use
the quasi-isomorphism \eqref{cpc:eq:operatorKoszul}, also after
reduction modulo $\ell$.
\end{proof}

\subsection{Proof of the scalar formula}

Let $k$ be an algebraic closure of $\mathbb F_\ell$. Since
$\ell\nmid\varphi(m)$, the character projectors for $U(m)$ are
defined over $k$. The conductor normal-form argument applies
verbatim in this characteristic: it only divides by unit-group
orders, and its edges do not divide by a weight or an Euler factor.
For clarity, the character $\chi$ has vectors
\[
 Y_{b,\chi}=\sum_{a\in U(b)}\chi(a)e_{a/b},
 \qquad \operatorname{cond}(\chi)\mid b\mid m.
\]
The prime edges are
$Y_{pb,\chi}-a_pY_{b,\chi}$ when $p\mid b$, and
$Y_{pb,\chi}-(a_p-\chi(p))Y_{b,\chi}$ when $p\nmid b$.
Successive elimination leaves exactly one conductor generator per
character. Therefore $D_m(k,\overline{\mathbf a})$ has one
one-dimensional summand for every character of $U(m)$, independently
of its weights. On that summand, $V_p$ acts by $\chi(p)$.

The complex in \eqref{cpc:eq:tateKoszul} consequently splits into
scalar Koszul complexes on
\begin{equation}\label{cpc:eq:characterparameters}
 (\chi(p)-\overline a_p:p\mid d).
\end{equation}
A nonzero scalar makes a Koszul complex contractible. If all these
scalars vanish, its differential is zero and its degree-$j$
homology has dimension $\binom rj$. Thus each surviving character
contributes $2^{r-1}$ to either cohomology parity.

The prescribed values in \eqref{cpc:eq:characterparameters} occur
if and only if \eqref{cpc:eq:compatibility} holds. A character of
the subgroup $H$ extends to $U(m)$ over $k$: one may successively
adjoin group elements and take the required roots in the algebraically
closed field. All relevant orders are prime to $\ell$. The set of
extensions is a torsor under the character group of $U(m)/H$, so
its cardinality is $\varphi(m)/|H|$. This proves the torsion
multiplicity in \eqref{cpc:eq:scalar}; extension to $k$ preserves
the dimension over $\mathbb F_\ell$.

For the free rank, rational coinvariants are an exact functor because
averaging over a finite group is possible over $\mathbb Q$. A
permutation set with $b$ points and $f$ fixed points has
$(b+(\ell-1)f)/\ell$ orbits. Its alternating count in the resolution
gives
\[
 \operatorname{rank}_{\mathbb Z}D/(T-1)D
 =\frac1\ell\sum_{S\subseteq P(q)}(-1)^{|S|}
 \left(\frac q{p_S}+(\ell-1)
           \frac m{p_{S\cap P(m)}}\right)
 =\frac{\varphi(q)}\ell.
\]
The second alternating sum vanishes because $r=\omega(d)>0$.
The classification of finitely generated abelian groups and
Lemma~\ref{cpc:lem:tate} complete the proof of
Theorem~\ref{cpc:thm:scalar}.

\subsection{All finite integral jets}

Let $B_L=\mathbb Z[t]/(t^L)$ with $L\ge1$, and take arbitrary
$a_p(t)\in B_L$. Test compatibility using the constant terms
$\overline a_p(0)$. If compatibility fails, put $c=0$ and
$\nu=0$. If it holds, let $c=\varphi(m)/|H|$ and define
\begin{equation}\label{cpc:eq:contact}
 \nu=\min_{p\mid d}
 \operatorname{ord}_t\bigl(\overline a_p(t)-\overline a_p(0)\bigr)
 \quad\text{in }\mathbb F_\ell[t]/(t^L),
 \qquad\operatorname{ord}_t(0)=L.
\end{equation}
In the compatible case $1\le\nu\le L$.

\begin{theorem}[Cyclic-prime integral jet formula]
\label{cpc:thm:jets}
Under \eqref{cpc:eq:hypotheses}, the underlying abelian group is
\begin{equation}\label{cpc:eq:jets}
 \boxed{\displaystyle
 D_q(B_L,\mathbf a)/(T-1)D_q(B_L,\mathbf a)
 \simeq\mathbb Z^{L\varphi(q)/\ell}
       \oplus(\mathbb Z/\ell\mathbb Z)^{\,c2^{r-1}\nu}.}
\end{equation}
More precisely, its integer torsion, as a $B_L$-module, is
\begin{equation}\label{cpc:eq:jetmodule}
 \left(B_L/(\ell,t^\nu)\right)^{c2^{r-1}}.
\end{equation}
The free summand in \eqref{cpc:eq:jets} is only asserted to be
free over $\mathbb Z$, not over $B_L$.
\end{theorem}
\begin{proof}
Over $k[t]/(t^L)$, the same conductor projectors split the fixed
module into free rank-one character blocks. A block whose constants
do not match has a unit parameter and is contractible. For each of
the $c$ matching characters, all parameters are the same list
\[
 -(\overline a_p(t)-\overline a_p(0)),\qquad p\mid d.
\]
In the principal ideal ring $k[t]/(t^L)$, invertible operations on
this list transform it into $(t^\nu,0,\ldots,0)$, interpreting
$t^L$ as zero. Both the kernel and the cokernel of multiplication
by $t^\nu$ are isomorphic to $k[t]/(t^\nu)$. Tensoring with the
$r-1$ remaining zero-differential exterior factors gives
\[
 H_j\simeq
 \left(k[t]/(t^\nu)\right)^{\binom rj}
\]
in each matching block. Lemma~\ref{cpc:lem:tate} gives the claimed
multiplicity.

For the $B_L$-module assertion no splitting-field ambiguity is
needed. The compatible character of $H$ takes values in
$\mathbb F_\ell$, so its averaging idempotent is defined over
$\mathbb F_\ell[t]/(t^L)$. Its image in the fixed distribution
module is a projective module of rank $c$ over this local ring,
hence is free of rank $c$. On that image every $V_p$ is the
prescribed scalar. The other character blocks are contractible
after the faithfully flat splitting extension and therefore have
zero homology already before that extension. The preceding scalar
Koszul calculation can consequently be performed over
$\mathbb F_\ell[t]/(t^L)$ itself, proving
\eqref{cpc:eq:jetmodule}. Finally the rational orbit count in the
scalar proof is multiplied by $L$, giving the free rank.
\end{proof}

\subsection{Explicit examples and boundaries of the result}

\paragraph{Level $91$, a genuinely odd torsion group.}
Take $q=7\cdot13$, $\ell=3$, and $u=16$. Then $u$ has order three
and $(u-1,91)=1$. For $a_7=a_{13}=1$,
\[
 D_{91}/(T-1)D_{91}\simeq\mathbb Z^{24}\oplus(\mathbb Z/3)^2.
\]
If either of these weights is not one modulo three, the torsion
vanishes. At jet length $L=4$, choose
$a_7(t)=1+t^2$ and $a_{13}(t)=1+3t+t^3$. Their contact order is
$\nu=2$, so the quotient is
$\mathbb Z^{96}\oplus(\mathbb Z/3)^4$.

\paragraph{Level $455$, a compatibility condition involving two primes.}
Take $q=5\cdot7\cdot13$, again with $\ell=3$ and $u=16$.
Here $m=5$, $d=91$, and
$H=\langle2,3\rangle=U(5)$ has order four.
Because $3=2^{-1}$ in $U(5)$ and every nonzero element of
$\mathbb F_3$ is its own inverse, compatibility is precisely
\[
 a_7\equiv a_{13}\not\equiv0\pmod3.
\]
Thus the quotient is
$\mathbb Z^{96}\oplus(\mathbb Z/3)^2$ when the two residues
are both one or both two, and is $\mathbb Z^{96}$ otherwise.
The value of $a_5$ is irrelevant. The case where both residues are
two shows why the criterion cannot be replaced by an independent
requirement $a_p\equiv1$ when a nontrivial fixed denominator remains.

\paragraph{The hypotheses exclude ramified local actions.}
At $q=9$, $\ell=3$, $u=4$, the fixed denominator is three, but it
is not relatively prime to its complementary factor. The five grid
orbits give generators $e_0,e_{1/3},e_{2/3},A,B$, where $A$ and
$B$ correspond to the two orbits of size three. The three raw
relations are
\[
 (1-a)e_0+e_{1/3}+e_{2/3}=0,\qquad
 3A-ae_{1/3}=0,\qquad 3B-ae_{2/3}=0.
\]
At $a=0,1,2$ these give, respectively,
\[
 \mathbb Z^2\oplus(\mathbb Z/3)^2,\qquad
 \mathbb Z^2\oplus\mathbb Z/3,\qquad
 \mathbb Z^2.
\]
These exact presentations explain why the fixed-grid hypotheses
must remain part of the statement. They are not counterexamples
to a claim in the inspected manuscript, which did not assert the
cyclic-prime formula.

\subsection{Verification, integration, and remaining questions}

The accompanying program constructs the complete defining rows
directly on $X_q$, identifies points in the orbits of $u$, and
computes the Smith normal form over $\mathbb Z$. It imports no
normal-form code from either incoming integral-distribution package.
Its frozen suite checks scalar weights for symmetry primes
$2,3,5$, prime-power and composite levels through $455$, incompatible
weight assignments, and nonconstant integral jets through length
four. The checks verify every nonzero invariant, not merely ranks
modulo $\ell$. They provide independent arithmetic regression
evidence; the all-level proofs are the arguments above.

The preceding exact results resolve the proposed prime-order subgroup
problem under \eqref{cpc:eq:hypotheses}. They do not settle it for an
arbitrary unit of prime order. Three concrete next problems remain.
First, compute the fixed complex for a local action of order $\ell$
on an $\ell$-power denominator, beginning with the displayed level-nine
presentation. Second, remove $\ell\nmid\varphi(m)$: the operator
Koszul model still applies, but its modular unit-group action need
not split into character lines. Third, compute the full module over
a multivariate deformation ring, where the scalar Koszul parameters
generate a nonprincipal ideal and a single contact order no longer
determines its homology. The first two problems concern new
stabilizer or modular representation structure; the third concerns
the syzygies of an explicit ideal.
